\documentclass[11pt]{beamer}

\usetheme{Madrid}
\usecolortheme{default}
\definecolor{unicaucaverde}{RGB}{0,90,60}
\setbeamercolor{structure}{fg=unicaucaverde}
\setbeamertemplate{navigation symbols}{}

\usepackage[utf8]{inputenc}
\usepackage[spanish]{babel}
\usepackage{amsmath,amssymb}

\newcommand{\R}{\mathbb{R}}
\newcommand{\vv}[1]{\mathbf{#1}}

\title[Semana 2: Gauss-Jordan]{Eliminación de Gauss-Jordan y sistemas homogéneos}
\subtitle{Álgebra Lineal para Ingeniería --- Semana 2 de 16}
\author{Universidad del Cauca}
\date{2026}

\begin{document}

\begin{frame}
  \titlepage
\end{frame}

\begin{frame}{Semana 2: dónde estamos}
  \begin{itemize}
    \item Unidad 1 de 8: Sistemas de ecuaciones lineales.
    \item Subtemas 1.2b, 1.3 y 1.4 --- RA1.2b, RA1.3, RA1.4.
    \item 4 horas: sistemas equivalentes/clasificación $\to$ Gauss-Jordan $\to$ sistemas homogéneos.
  \end{itemize}
  \tableofcontents
\end{frame}

\section{Sistemas equivalentes y clasificación}

\begin{frame}{Sistemas equivalentes}
  \begin{block}{Definición}
    Dos sistemas son \textbf{equivalentes} si tienen el mismo conjunto de soluciones.
  \end{block}
  \vfill
  \begin{alertblock}{Teorema}
    Cada operación elemental de fila produce un sistema equivalente:
    \begin{enumerate}
      \item Intercambiar dos filas.
      \item Multiplicar una fila por $c\neq 0$.
      \item Sumar a una fila un múltiplo de otra.
    \end{enumerate}
  \end{alertblock}
  \vfill
  $\Rightarrow$ Al final de Gauss-Jordan el sistema resultante es \textbf{equivalente} al original.
\end{frame}

\begin{frame}{Clasificación: consistente e inconsistente}
  \begin{block}{Sistema consistente}
    Tiene \textbf{al menos una solución} (puede ser única o infinitas).
  \end{block}
  \begin{alertblock}{Sistema inconsistente}
    \textbf{No tiene ninguna solución.} Se detecta en la forma escalonada:
    aparece una fila $[\,0\;\cdots\;0\mid c\,]$ con $c\neq 0$ (la ecuación $0=c$).
  \end{alertblock}
  \vfill
  \textbf{Ejemplo inconsistente:}
  \[
    x_1+x_2=3,\quad 2x_1+2x_2=7
    \;\xrightarrow{F_2-2F_1}\;
    \left(\begin{array}{cc|c}1&1&3\\0&0&1\end{array}\right)
  \]
  La segunda fila dice $0=1$: contradicción.
\end{frame}

\begin{frame}{Ejemplo 1 — Sistema inconsistente (solución)}
  \[
    x_1+x_2=3, \qquad 2x_1+2x_2=7
  \]
  \[
    \left(\begin{array}{cc|c}1&1&3\\2&2&7\end{array}\right)
    \xrightarrow{F_2-2F_1}
    \left(\begin{array}{cc|c}1&1&3\\0&0&1\end{array}\right)
  \]
  \begin{alertblock}{Resultado}
    Fila $[0\ 0\mid 1]$: ecuación $0=1$ \textbf{(contradicción).}\\
    El sistema es \textbf{inconsistente}: \textbf{no tiene solución}.\\
    Geométricamente: las dos rectas son \textbf{paralelas}.
  \end{alertblock}
\end{frame}

\begin{frame}{Ejemplo 2 — Solución única}
  \[
    x_1+x_2+x_3=6,\quad 2x_1-x_2+x_3=3,\quad x_1+2x_2-x_3=2
  \]
  Aplicando Gauss-Jordan (ver \texttt{notas.tex} Ejemplo 2):
  \[
    \cdots\longrightarrow
    \left(\begin{array}{ccc|c}1&0&0&1\\0&1&0&2\\0&0&1&3\end{array}\right)
  \]
  \begin{block}{Resultado}
    Rango $r=3=n$: sistema \textbf{consistente, solución única}:
    \[
      x_1=1,\quad x_2=2,\quad x_3=3.
    \]
  \end{block}
\end{frame}

\begin{frame}{Ejemplo 3 — Infinitas soluciones}
  \[
    x_1+2x_2-x_3=0, \qquad 2x_1+4x_2-2x_3=0
  \]
  \[
    \left(\begin{array}{ccc|c}1&2&-1&0\\2&4&-2&0\end{array}\right)
    \xrightarrow{F_2-2F_1}
    \left(\begin{array}{ccc|c}1&2&-1&0\\0&0&0&0\end{array}\right)
  \]
  \begin{block}{Resultado}
    Rango $r=1<n=3$: \textbf{2 variables libres} ($x_2=s$, $x_3=t$).\\
    Sistema \textbf{consistente, infinitas soluciones}:
    \[
      \{(-2s+t,\;s,\;t)\mid s,t\in\mathbb{R}\}.
    \]
  \end{block}
  \small Ej.: $(0,0,0)$ con $s=t=0$;\quad $(1,0,1)$ con $s=0,\,t=1$.
\end{frame}

\section{Gauss-Jordan}


\begin{frame}{Matriz aumentada}
  \[
    [A\mid\vv{b}] =
    \left(\begin{array}{ccc|c}
      a_{11} & \cdots & a_{1n} & b_1 \\
      \vdots & \ddots & \vdots & \vdots \\
      a_{m1} & \cdots & a_{mn} & b_m
    \end{array}\right)
  \]
  Contiene toda la información del sistema, sin las incógnitas.
\end{frame}

\begin{frame}{Operaciones elementales de fila}
  \begin{enumerate}
    \item Intercambiar dos filas.
    \item Multiplicar una fila por un escalar $\neq 0$.
    \item Sumar a una fila un múltiplo de otra.
  \end{enumerate}
  \vfill
  Ninguna cambia el conjunto solución del sistema.
\end{frame}

\begin{frame}{Forma escalonada reducida (RREF)}
  \begin{itemize}
    \item Pivote de cada fila no nula = 1.
    \item Cada pivote a la derecha del anterior.
    \item Filas nulas al final.
    \item Columna de cada pivote: ceros en las demás posiciones.
  \end{itemize}
  \begin{alertblock}{Gauss-Jordan}
    Algoritmo que lleva $[A\mid\vv{b}]$ a su RREF mediante operaciones elementales.
  \end{alertblock}
\end{frame}

\begin{frame}{Número de soluciones}
  Sea $r$ = pivotes de $[A\mid\vv{b}]$, $r'$ = pivotes de $A$.
  \begin{itemize}
    \item $r' < r$ $\Rightarrow$ \textbf{inconsistente}.
    \item consistente y $r=n$ $\Rightarrow$ \textbf{solución única}.
    \item consistente y $r<n$ $\Rightarrow$ \textbf{infinitas} ($n-r$ variables libres).
  \end{itemize}
\end{frame}

\begin{frame}{Ejemplo}
  \[
  \begin{aligned}
    x_1+x_2+x_3 &= 6 \\
    2x_1-x_2+x_3 &= 3 \\
    x_1+2x_2-x_3 &= 2
  \end{aligned}
  \]
  Reduciendo la matriz aumentada se llega a solución única $(1,2,3)$.
  \vfill
  \textbf{Taller (hora 2):} \texttt{visualizacion/} sección \emph{1.3 Gauss-Jordan} --- reducir
  paso a paso.
\end{frame}

\begin{frame}{Aplicación: corrientes de malla en un circuito}
  Ley de voltajes de Kirchhoff en las tres mallas $\Rightarrow$ sistema en $I_1,I_2,I_3$:
  \[
  \begin{aligned}
    2I_1 - I_2 &= 5 \\
    -I_1 + 3I_2 - I_3 &= 0 \\
    -I_2 + 2I_3 &= 1
  \end{aligned}
  \]
  \begin{alertblock}{Sin método sistemático}
    Sustitución: funciona, pero exige decidir un orden de variables --- se complica al crecer
    el circuito.
  \end{alertblock}
\end{frame}

\begin{frame}{Aplicación: misma solución por Gauss-Jordan}
  Reduciendo $[A\mid\vv{b}]$ con \textbf{fracciones exactas} (sin redondear a decimal en ningún
  paso) se llega a
  \[
    I_1=\frac{13}{4}, \qquad I_2=\frac{3}{2}, \qquad I_3=\frac{5}{4} \quad \text{(A)}
  \]
  igual que por sustitución --- pero de forma sistemática, sin depender del orden elegido.
  \vfill
  \textbf{Taller:} \texttt{visualizacion/} sección \emph{1.3 Gauss-Jordan} --- cargar el ejemplo
  del circuito y reducir paso a paso.
\end{frame}

\section{Sistemas homogéneos}

\begin{frame}{Sistema homogéneo: \texorpdfstring{$A\vv{x}=\vv{0}$}{Ax=0}}
  \[
  \begin{aligned}
    a_{11}x_1+\cdots+a_{1n}x_n &= 0 \\
    &\;\vdots \\
    a_{m1}x_1+\cdots+a_{mn}x_n &= 0
  \end{aligned}
  \]
  \begin{alertblock}{Solución trivial}
    $\vv{x}=\vv{0}$ siempre es solución: todo sistema homogéneo es consistente.
  \end{alertblock}
  La pregunta relevante: ¿hay soluciones \textbf{no triviales}?
\end{frame}

\begin{frame}{¿Cuándo hay soluciones no triviales?}
  \begin{alertblock}{Teorema}
    Hay soluciones no triviales $\iff$ $r<n$ (variables libres). Si $m<n$, siempre las hay.
  \end{alertblock}
  \vfill
  \textbf{Ejemplo:}
  \[
    x_1+2x_2-x_3=0, \qquad 2x_1+4x_2-2x_3=0
  \]
  (la segunda es el doble de la primera) $\Rightarrow$ $r=1<n=3$: infinitas soluciones,
  parametrizadas por $x_2=s$, $x_3=t$: $(-2s+t,\,s,\,t)$.
\end{frame}

\section{Soluciones de ejercicios}

\begin{frame}{Ej. 1 — ¿Inconsistente o consistente?}
  $2x_1-x_2+x_3=4,\quad x_1+3x_3=-1,\quad -x_1+x_2+2x_3=0$
  \[
    \left(\begin{array}{ccc|c}1&0&3&-1\\2&-1&1&4\\-1&1&2&0\end{array}\right)
    \xrightarrow{F_2-2F_1,\;F_3+F_1}
    \left(\begin{array}{ccc|c}1&0&3&-1\\0&-1&-5&6\\0&1&5&-1\end{array}\right)
    \xrightarrow{F_3+F_2}
    \left(\begin{array}{ccc|c}1&0&3&-1\\0&-1&-5&6\\0&0&0&5\end{array}\right)
  \]
  \begin{alertblock}{Resultado}
    Fila $[0\ 0\ 0\mid 5]$: ecuación $0=5$. \textbf{Sistema inconsistente — sin solución.}
  \end{alertblock}
\end{frame}

\begin{frame}{Ej. 2 — Solución única}
  $x+y+z=2,\quad 2x-y+z=1,\quad x+2y-z=3$
  \[
    \cdots\longrightarrow
    \left(\begin{array}{ccc|c}1&1&1&2\\0&1&-2&1\\0&0&1&0\end{array}\right)
    \longrightarrow
    \left(\begin{array}{ccc|c}1&0&0&1\\0&1&0&1\\0&0&1&0\end{array}\right)
  \]
  \begin{block}{Resultado}
    \[ \boxed{x=1,\quad y=1,\quad z=0.} \]
  \end{block}
  \small(Ver \texttt{notas.tex} §\,Soluciones, Ej.\ 2 para el paso completo.)
\end{frame}

\begin{frame}{Ej. 3 — Tres casos según $k$}
  $x+ky=1,\quad kx+y=1$
  \[
    \left(\begin{array}{cc|c}1&k&1\\k&1&1\end{array}\right)
    \xrightarrow{F_2-kF_1}
    \left(\begin{array}{cc|c}1&k&1\\0&1-k^2&1-k\end{array}\right)
  \]
  \begin{itemize}
    \item $k=1$: fila 2 $\to [0\ 0\mid 0]$ \textbf{→ infinitas} ($x+y=1$).
    \item $k=-1$: fila 2 $\to [0\ 0\mid 2]$ \textbf{→ inconsistente}.
    \item $k\neq\pm 1$: \textbf{solución única} $x=y=\dfrac{1}{1+k}$.
  \end{itemize}
\end{frame}

\begin{frame}{Ej. 4 — Justificación sin resolver}
  \textbf{Sistema homogéneo}: $m=3$ ecuaciones, $n=5$ incógnitas.
  \vfill
  \begin{block}{Argumento}
    El rango máximo es $r\leq m=3 < 5=n$.\\
    Por tanto $n-r\geq 2$ \textbf{variables libres} → infinitas soluciones no triviales.\\
    \textbf{No hace falta reducir la matriz para saberlo.}
  \end{block}
\end{frame}

\begin{frame}{Ej. 5 — Infinitas soluciones (homogéneo)}
  $x_1-x_2+2x_3=0,\quad 2x_1-2x_2+4x_3=0$
  \[
    \left(\begin{array}{ccc|c}1&-1&2&0\\2&-2&4&0\end{array}\right)
    \xrightarrow{F_2-2F_1}
    \left(\begin{array}{ccc|c}1&-1&2&0\\0&0&0&0\end{array}\right)
  \]
  \begin{block}{Resultado — variables libres: $x_2=s$, $x_3=t$}
    \[ x_1 = s-2t \]
    \[\boxed{\{(s-2t,\;s,\;t)\mid s,t\in\mathbb{R}\}}\]
  \end{block}
  \small Ej.: $(0,0,0)$ trivial;\quad $(1,1,0)$ no trivial;\quad $(-2,0,1)$ no trivial.
\end{frame}

\begin{frame}{Ej. 6 — ¿Cuántos pivotes?}
  Sistema homogéneo $3\times 3$ con \textbf{única solución}: la trivial $\vv{x}=\vv{0}$.
  \vfill
  \begin{block}{Argumento}
    Solución única $\Rightarrow$ ninguna variable libre $\Rightarrow$ ninguna columna sin
    pivote.\\
    $n=3$ incógnitas, 0 variables libres $\Rightarrow$ \textbf{$r=3$ pivotes}.\\
    Equivalentemente: la matriz $A$ es cuadrada no singular ($\det A\neq 0$).
  \end{block}
\end{frame}

\begin{frame}{Cierre de la Unidad 1}
  \begin{itemize}
    \item RA1.1--RA1.2 (Semana 1): forma general, sistemas $2\times 2$, tres casos geométricos.
    \item RA1.2b (Semana 2): sistemas equivalentes, consistente/inconsistente.
    \item RA1.3--RA1.4 (Semana 2): Gauss-Jordan, número de soluciones, sistemas homogéneos.
    \item Quiz de la Unidad 1 la próxima sesión.
  \end{itemize}
  \vfill
  \centering \Large Unidad 2: Vectores y matrices.
\end{frame}

\end{document}
